\begin{figure}[!ht]
\centering
\begin{tikzpicture}[
  every node/.style={font=\small},
  field/.style={inner sep=4pt},
  inclusion/.style={-{Stealth[length=4pt]},semithick}
]
\begin{scope}[xshift=-3.45cm]
  \node[font=\small\bfseries] at (0,4.65) {Characteristic zero};
  \node[field] (kb) at (0,0) {$K_b$};
  \node[field] (ks) at (-1.45,1.35) {$K_s$};
  \node[field] (kbp) at (1.45,1.35) {$K_b'$};
  \node[field] (l) at (-1.45,2.8) {$L$};
  \node[field] (lp) at (0,4) {$L'=LK_b'$};
  \draw[inclusion] (kb)--(ks);
  \draw[inclusion] (kb)--(kbp);
  \node[font=\footnotesize,align=left] at (1.85,.35)
    {unramified\\degree $n$};
  \draw[inclusion] (ks)--(l);
  \draw[inclusion] (l)--(lp);
  \draw[inclusion] (kbp)--(lp);
\end{scope}
\begin{scope}[xshift=3.45cm]
  \node[font=\small\bfseries] at (0,4.65) {Residue fields};
  \node[field] (f0) at (0,0) {$F_0$};
  \node[field] (fp) at (1.45,1.35) {$F'$};
  \node[field] (f) at (-1.45,2.8) {$F$};
  \node[field] (s) at (0,4) {$S=FF'$};
  \draw[inclusion] (f0)--(fp);
  \draw[inclusion] (f0)--(f);
  \draw[inclusion] (f)--(s);
  \draw[inclusion] (fp)--(s);
\end{scope}
\node[font=\footnotesize,align=center] at (0,-.85)
 {$\operatorname{res}(K_b)=\operatorname{res}(K_s)=F_0,
   \qquad [F:F_0]=[L:K_s],\qquad [L:K_s]_p<p^h.$};
\end{tikzpicture}
\caption{The auxiliary fields in a descent arising from a small Sylow orbit. Arrows denote field inclusions. The right-hand diagram records the residues of the left-hand fields; in particular $K_b$ and $K_s$ have the same residue. The field $K_b'$ splits the marked sections, and its residue $F'$ splits their reductions. The cyclic extension $K_b'/K_b$ reduces to $F'/F_0$, and the Galois actions of $L'/L$ and $S/F$ induce the same permutations of the marked points.}
\label{fig:field-tower}
\end{figure}
